\section*{Chapitre \ref{ch:qm:multivariate-series-and-cointegration} --- Séries multivariées et cointégration}

\begin{solution}{exo:qm:multivariate-series-and-cointegration:1}
Les valeurs propres de la première matrice sont $0.9$ et $0.1$~: stationnaire. Celles de la seconde valent $1.0$ et $0.2$~: une \omterm{def:qm:linear-time-series:unitroot}{racine unitaire}, donc non stationnaire~; la somme $x + y$
suit une marche aléatoire tandis que la différence $x - y$ est stationnaire (de coefficient 0.2)~: les deux séries sont cointégrées avec $\beta = (1, -1)$.
\end{solution}

\begin{solution}{exo:qm:multivariate-series-and-cointegration:2}
$k = 2$, $r = 1$ (le vecteur $(1, -0.9)$ à un facteur d'échelle près), et $k - r = 1$ tendance commune.
\end{solution}

\begin{solution}{exo:qm:multivariate-series-and-cointegration:3}
$-3.74066 - 8.5631/500 - 10.852/500^2 + 27.982/500^3 = -3.758$.
\end{solution}

\begin{solution}{exo:qm:multivariate-series-and-cointegration:4}
$\Delta y_t = -0.5(y_{t-1} - 0.9x_{t-1}) + e_t$ et $\Delta x_t = u_t$. Ainsi $\beta = (1, -0.9)^\top$ et $\alpha = (-0.5, 0)^\top$ pour $(y, x)$~: $y$ corrige la moitié de tout écart
à chaque période~; $x$ ne corrige pas du tout (il est faiblement exogène).
\end{solution}

\begin{solution}{exo:qm:multivariate-series-and-cointegration:5}
La factorisation de Cholesky écrit le second résidu comme $0.9z_1 + \sqrt{0.19}\,z_2$, de sorte que 81~\% de sa variance à un pas soit attribuée au premier choc~; avec l'ordre inverse
elle serait de 19~\% (ou 0~\% de la part de la première variable attribuée au second choc dans l'ordre d'origine).
\end{solution}

\begin{solution}{exo:qm:multivariate-series-and-cointegration:6}
Statistique de la trace pour $r = 0$~: $-1\,000(\ln0.95 + \ln0.99 + \ln0.998) = 63.35$~; $r \le 1$~: 12.05~; $r \le 2$~: 2.00. Face à 35.23, 20.36 et 9.23~: on rejette $r = 0$ mais pas $r \le 1$, donc rang un.
\end{solution}

\begin{solution}{exo:qm:multivariate-series-and-cointegration:7}
\texttt{rolling\_weights()}~: aucune \omterm{def:qm:multivariate-series-and-cointegration:coint}{cointégration} à 5~\% dans 20 des 46 fenêtres~; lorsqu'un vecteur est trouvé, le poids à 2 ans varie de 0.27 à 0.63 et le poids à 10 ans
de 0.40 à 0.97.
\end{solution}

\begin{solution}{exo:qm:multivariate-series-and-cointegration:8}
Le résidu d'une régression estimée doit être jugé à l'aide des valeurs critiques d'Engle--Granger, $-3.74$ à 5~\% pour trois variables, et non $-2.86$~; $-3.1$ ne rejette pas
l'\omterm{def:qm:testing-and-multiple-testing:test}{hypothèse nulle} d'absence de \omterm{def:qm:multivariate-series-and-cointegration:coint}{cointégration}. Avec la valeur de Dickey--Fuller, des marches aléatoires indépendantes franchiraient le seuil dans 27~\% des cas. Il faut aussi vérifier la stabilité des poids avant
de négocier le butterfly.
\end{solution}

\begin{solution}{pb:qm:multivariate-series-and-cointegration:1}
\textbf{1.} Statistiques de Dickey--Fuller $-1.30$, $-1.33$, $-1.31$~: aucun rendement ne rejette l'hypothèse de \omterm{def:qm:linear-time-series:unitroot}{racine unitaire}.
\textbf{2.} Le critère AIC continue de s'améliorer jusqu'au maximum de dix retards~; le VAR(10) des variations est stable.
\textbf{3.} 7.9, 7.6 et 7.1 pb~; corrélations 0.90 (2 ans--5 ans), 0.82 (2 ans--10 ans), 0.94 (5 ans--10 ans).
\textbf{4.} Avec le rendement à 2 ans ordonné en premier, un choc de 7.9 pb déplace cumulativement les rendements à 2, 5 et 10 ans de 10.1, 8.3 et 6.7 pb sur dix jours, et explique 67~\% de la
variance des erreurs de prévision à dix jours du rendement à 10 ans.
\textbf{5.} Les six directions sont significatives à 5~\% ($F$ de 1.88 à 7.16 avec 10 et 12\,532 degrés de liberté)~: avec tant de données, de faibles effets d'avance-retard issus de décalages
de dates et de nouvelles communes sont détectables~; il s'agit de précédence prédictive, non de causalité.
\textbf{6.} 59.77, 14.93 et 2.34 face à 35.23, 20.36 et 9.23~: rang un.
\textbf{7.} $0.41 \times$ 2 ans $- 1 \times$ 5 ans $+ 0.61 \times$ 10 ans~: les ailes totalisent 1.02 (neutre au niveau), et l'inclinaison vers le 10 ans neutralise aussi l'essentiel de la
pente.
\textbf{8.} L'échéance 5 ans sur les autres~: $0.40$ et $0.62$, statistique du résidu $-10.04$ face à une valeur critique à 5~\% de $-3.74$.
\textbf{9.} 43 jours ouvrés (coefficient AR 0.9841) face à 63 (0.9891).
\textbf{10.} 2.06 et 2.15 pb par jour (le butterfly 1-2-1 ramené au même poids de l'échéance 5 ans).
\textbf{11.} Dans 26 des 46 fenêtres (20 n'en trouvent aucune à 5~\%).
\textbf{12.} Le poids à 2 ans varie de 0.27 à 0.63 et celui à 10 ans de 0.40 à 0.97, parmi les fenêtres qui en trouvent un.
\textbf{13.} Les moindres carrés choisissent la combinaison qui paraît la plus stationnaire dans l'échantillon, si bien que la statistique du résidu est repoussée vers la gauche même lorsque rien n'est
cointégré~; plus il y a de variables, plus le décalage est grand.
\textbf{14.} Appliquer la correction de Kendall $(1 + 3\hat\rho)/n$ au coefficient AR(1) du butterfly~: 0.9841 devient 0.9844 avec 12\,574 jours et la \omterm{def:qm:stochastic-differential-equations:ou}{demi-vie} 44 jours
au lieu de 43~; puis simuler sa distribution comme au chapitre~17.
\textbf{15.} Que les poids des années suivantes diffèrent, de sorte que le butterfly négocié porte une exposition au niveau ou à la pente qu'il devait éliminer, et revienne à la moyenne plus lentement que le
backtest l'annonçait.
\textbf{16.} Le butterfly de Johansen, dont les poids sont réestimés sur une fenêtre récente et bornés près des valeurs de long terme, recoupé avec des poids neutres en duration.
\textbf{17.} Par l'écart-type stationnaire du butterfly, $2.06/\sqrt{1 - 0.9841^2} = 11.6$ pb~: dimensionner la position de sorte qu'un mouvement de deux ou trois écarts-types
(25 à 35 pb) soit supportable, et s'attendre à ce qu'un écart typique se réduise de moitié en deux mois environ.
\textbf{18.} Parce que les charges factorielles de la courbe changent selon le régime de taux~: réestimer sur une fenêtre glissante chaque mois ou chaque trimestre, et ne négocier que lorsque le test de rang
continue de détecter le vecteur.
\textbf{19.} Résultat à retenir~: \emph{le butterfly qui revient à la moyenne}~: le butterfly de \omterm{def:qm:multivariate-series-and-cointegration:coint}{cointégration} de Johansen des rendements à 2, 5 et 10 ans, 1976--2026, pondère les ailes de $0.41$ et $0.61$ par
unité de l'échéance 5 ans et revient à la moyenne avec une \omterm{def:qm:stochastic-differential-equations:ou}{demi-vie} de 43 jours ouvrés, contre 63 pour le butterfly 1-2-1.
\textbf{20.} Il promet que les écarts d'un panier sont temporaires en moyenne~; il ne promet ni combien de temps ils durent, ni l'ampleur qu'ils atteindront d'abord, ni que les poids resteront
les mêmes.
\end{solution}

\begin{solution}{iq:qm:multivariate-series-and-cointegration:1}
La corrélation décrit le comouvement des variations sur un horizon~; la \omterm{def:qm:multivariate-series-and-cointegration:coint}{cointégration} signifie que les niveaux partagent des tendances stochastiques, de sorte qu'une combinaison d'entre eux est stationnaire. Deux
séries peuvent être fortement corrélées et s'écarter l'une de l'autre indéfiniment, ou faiblement corrélées au jour le jour et liées entre elles à long terme.
\iqlookfor{Variations contre niveaux, et un contre-exemple dans les deux sens.}
\end{solution}

\begin{solution}{iq:qm:multivariate-series-and-cointegration:2}
Choisir des poids qui rendent le butterfly stationnaire et à retour rapide vers la moyenne~: estimer le \omterm{def:qm:multivariate-series-and-cointegration:coint}{vecteur de cointégration} (Johansen ou Engle--Granger avec les bonnes valeurs critiques),
comparer avec des poids neutres en duration ou neutres en ACP, vérifier la stabilité sur fenêtres glissantes, et mesurer la \omterm{def:qm:stochastic-differential-equations:ou}{demi-vie} et la volatilité du butterfly qui en résulte.
\iqlookfor{La stationnarité comme critère, le bon test, et la stabilité.}
\end{solution}

\begin{solution}{iq:qm:multivariate-series-and-cointegration:3}
Les coefficients ont été estimés pour rendre le résidu aussi petit, et donc d'apparence aussi stationnaire, que possible~; sous l'\omterm{def:qm:testing-and-multiple-testing:test}{hypothèse nulle}, la statistique du résidu est décalée vers la gauche, et le
décalage croît avec le nombre de régresseurs. Utiliser les valeurs critiques d'Engle--Granger (MacKinnon)~: $-3.34$ pour deux variables, $-3.74$ pour trois à 5~\%, asymptotiquement.
\iqlookfor{L'argument du pré-ajustement et la bonne table de valeurs critiques.}
\end{solution}

\begin{solution}{iq:qm:multivariate-series-and-cointegration:4}
Le fait que le passé d'une série améliore les prévisions d'une autre, compte tenu de son propre passé~; il teste une précédence prédictive dans l'ensemble d'information utilisé. Il n'établit pas
causalité~: des facteurs communs omis, des décalages de dates dans les données et des anticipations peuvent tous produire cet effet.
\iqlookfor{Prédiction et non causalité, et la mise en garde sur les variables omises.}
\end{solution}

\begin{solution}{iq:qm:multivariate-series-and-cointegration:5}
Parce que les chocs corrélés sont répartis par une factorisation de Cholesky~: la première variable reçoit toute la part commune du choc. Des ordres différents attribuent autrement la
corrélation contemporaine~; indiquer l'ordre retenu, essayer d'autres ordres, ou utiliser des réponses généralisées.
\iqlookfor{L'orthogonalisation et son arbitraire.}
\end{solution}

\begin{solution}{iq:qm:multivariate-series-and-cointegration:6}
Après avoir retiré les retards de court terme, calculer les corrélations canoniques entre les différences $\Delta Y_t$ et les niveaux retardés $Y_{t-1}$. Des combinaisons
stationnaires des niveaux prédisent les différences, elles ont donc de grandes corrélations canoniques~; le nombre de ces corrélations significativement non nulles est le rang, les
vecteurs propres estiment $\beta$, et la statistique de la trace somme $-T\ln(1 - \lambda_i)$ sur les plus petites.
\iqlookfor{Régression à rang réduit par corrélations canoniques.}
\end{solution}
